\section{From option pricing to physical commodity forecasts}
\label{sec:physical-measure-bridge}

The option evidence above identifies a pricing layer. The operating chapters
that follow require a different object: a distribution of future selling
prices from which revenue can be forecast. This section places the change of
measure at that interface, before production, costs and discounted cash flow.
It applies to all four model families and to both commodities. It is not an
editorial qualification appended to a completed corporate valuation.

The analysis in this section was updated on 6 October 2026 and keeps the
market cutoff of 2 September 2026. Joint identification, numerical
convergence of the richer models, commodity basis uncertainty and the
corporate cash-flow bridge remain open, so none of these diagnostics replaces
the scenario outputs with a validated corporate valuation.

\subsection{The economic object and the two probability measures}

An option calibration estimates coefficients under a pricing measure
$\mathbb Q$, conditional on the chosen model and market inputs. Expected
realised commodity prices and revenue forecasts require the physical measure
$\mathbb P$. An option-implied distribution alone does not identify expected
physical returns, volatility risk premia or jump risk premia. Joint
underlying and option evidence is therefore relevant to identification
\cite{pan2002}. The model families themselves are not restricted to
$\mathbb Q$: GBM, stochastic variance, Poisson jumps and Hawkes arrivals can
be specified under $\mathbb P$. What cannot be assumed is that their fitted
coefficients are identical under both measures.

There are two coherent valuation routes. A corporate DCF forecasts physical
cash flows and uses a discounting convention appropriate to their risk. A
pricing approach values appropriately specified claims under a pricing
measure and numeraire, or with a compatible stochastic discount factor.
Using option-implied paths and subsequently applying WACC does not perform
a change of measure. Nor does assigning independent stochastic WACC draws
to physical paths automatically identify a valid pricing kernel. The
cash-flow definition, risk adjustment and treatment of dependence must be
reconciled in either route.

The intended route here is physical commodity paths, realised selling-price
basis, attributable sales volumes, revenue, reconciled FCFF and corporate
discounting. Gold is required in USD per troy ounce and copper in USD per
pound, with explicit conversion to the financial reporting scale. GLD shares,
a copper futures delivery and Barrick's quarterly average realised selling
price are distinct economic objects. Restarting the first at the level of
the last does not resolve their differences.

\subsection{A restricted coefficient map for all four families}

The implemented starting specification preserves the same underlying and
event support across measures. Write the uncompensated pricing dynamics as
\begin{equation}
\frac{\dd S_t}{S_{t^-}}=(a_t^{Q}-\lambda_t^{Q}k_Q)\dd t
 +\sqrt{v_t}\dd W_t^{S,Q}+(e^J-1)\dd N_t,
\qquad k_Q=\mathbb E^Q[e^J-1].
\label{eq:pq-price-start}
\end{equation}
Here $a_t^Q$ is a supported carry specification; it is not WACC. For the
usual ETF formulation it is $r_t-q_t$. A fixed-delivery futures formulation
has zero pricing drift under deterministic rates; the slope across delivery
dates is not its physical time-series drift. Commodity spot carry may also
require storage costs and convenience yield \cite{schwartz1997}.

Use the Brownian convention
$\dd W_t^{S,Q}=\dd W_t^{S,P}+\eta_S\sqrt{v_t}\dd t$ and
$\dd W_t^{v,Q}=\dd W_t^{v,P}+\eta_v\sqrt{v_t}\dd t$.
For $|\rho|<1$, the two shifts are represented through the correlated
Brownian covariance matrix; they are not independent Brownian drivers when
$\rho\ne0$. Quadratic variation preserves $\xi$, $\rho$ and the diffusion
coefficient on the same path. Substitution into CIR variance gives
\begin{equation}
\kappa_P=\kappa_Q-\xi\eta_v>0,\qquad
\theta_P=\frac{\kappa_Q\theta_Q}{\kappa_P}.
\label{eq:pq-variance-map}
\end{equation}
Thus $\kappa\theta$ is preserved in this restricted family. This is a
consequence of the selected square-root shift, not a universal restriction
on every possible stochastic-volatility measure change. The estimation
parameterisation uses $\log(\kappa_P/\kappa_Q)$ to enforce positivity.

For normally distributed log jump sizes, a separate positive intensity
ratio $c_I$ and an exponential mark tilt $b_J$ specify
\begin{align}
\nu_P(\dd j)&=c_I\frac{e^{b_Jj}}{M_Q(b_J)}\nu_Q(\dd j),&
 M_Q(b_J)&=e^{b_J\mu_{J,Q}+b_J^2\sigma_J^2/2},\\
\lambda_t^P&=c_I\lambda_t^Q,&
\mu_{J,P}&=\mu_{J,Q}+b_J\sigma_J^2.
\label{eq:pq-jump-map}
\end{align}
The mark variance is unchanged under this tilt. Let
$k_P=\exp(\mu_{J,P}+\sigma_J^2/2)-1$. The physical expected instantaneous
price return, including jump compensation, is
\begin{equation}
\mu_t^P=a_t^Q+\eta_Sv_t+
 \lambda_t^Q(c_Ik_P-k_Q).
\label{eq:pq-total-drift}
\end{equation}
Changing only the diffusion drift while retaining an inconsistent jump
compensator would give the wrong physical mean and hence the wrong expected
revenue. Equivalently, the uncompensated physical drift is
$a_t^Q+\eta_Sv_t-\lambda_t^Qk_Q$, with arrivals and marks generated under
$\mathbb P$.

For exponential Hawkes arrivals, $\lambda_t^Q$ denotes the pricing-intensity
functional evaluated on the same realised counting history. Constant scaling
then preserves the exponential form:
\begin{equation}
\lambda_{0,P}=c_I\lambda_{0,Q},\quad
\bar\lambda_P=c_I\bar\lambda_Q,\quad
\alpha_P=c_I\alpha_Q,\quad\beta_P=\beta_Q.
\label{eq:pq-hawkes-map}
\end{equation}
Both branching ratios are constrained below one; in particular
$c_I\alpha_Q/\beta_Q<1$. Stationarity and the algebra above do not, by
themselves, prove that the candidate density is a true martingale over the
required horizon. The finite-horizon argument below establishes the density for the stated
restricted coefficient family. Initial-state consistency and empirical
validation remain separate requirements.

The Brownian part is checked separately using Novikov's sufficient condition.
With $C=\left(\begin{smallmatrix}1&\rho\\\rho&1\end{smallmatrix}\right)$,
set $\gamma=\tfrac12\eta^\top C^{-1}\eta$. The required CIR exponential
moment has Riccati equation
$B'=\gamma-\kappa_QB+\tfrac12\xi^2B^2$, $B(0)=0$.
If $\kappa_Q^2-2\xi^2\gamma\geq0$, this moment has no finite positive-time
explosion. Otherwise, with $d=\sqrt{2\xi^2\gamma-\kappa_Q^2}$, its explosion
time is $2[\pi/2+\arctan(\kappa_Q/d)]/d$; the tested horizon must lie strictly
below it. The three current conditional premium estimates pass the stronger
non-explosion criterion, including the five-year horizon. This establishes
the sufficient diffusion-density check for those coefficients; it does not
establish shared filtered states, empirical identification or estimation
convergence. The joint density requires the additional argument below.

\paragraph{Finite-horizon joint density.}
For the exact restricted family, the marked Girsanov formula
\cite{schonbucher2001} can be combined with a localization and entropy
argument. The need for a true density, rather than an algebraic intensity
change, is emphasized by Sokol and Hansen \cite{sokol-hansen2015}; the
following bounds are derived here for our specification. Assume finite
constant coefficients, deterministic finite initial states, nonnegative CIR
variance, unmarked linear Hawkes excitation, normal conditional marks and
$c_I>0$. Let $Z$ be the product of the Brownian stochastic exponential and
\begin{equation}
 Z_t^{J}=\exp\!\left\{
 \sum_{s\leq t}\left[\log c_I+b_JJ_s-b_J\mu_{J,Q}
                   -\tfrac12b_J^2\sigma_J^2\right]
 -(c_I-1)\int_0^t\lambda_s^Q\,\dd s\right\}.
\label{eq:pq-jump-density}
\end{equation}
Continuous and jump density local martingales have zero mutual quadratic
covariation; independence of their density factors is not required.
Take localizing times $\tau_n$ bounding the states and event counts and let
$\mathbb P_n$ have density $Z_{T\wedge\tau_n}$. With $\kappa_P>0$, dropping
the negative CIR mean-reversion term gives
$\mathbb E^{\mathbb P_n}[v_{t\wedge\tau_n}]
 \leq v_0+\kappa_Q\theta_Qt$.
Dropping the negative intensity decay term and applying Gronwall gives
$\mathbb E^{\mathbb P_n}[\lambda^Q_{t\wedge\tau_n}]
 \leq(\lambda_0^Q+\beta_Q\bar\lambda_QT)e^{c_I\alpha_Qt}$.
Hence, uniformly in $n$,
\begin{align}
 B_v(T)&=v_0T+\tfrac12\kappa_Q\theta_QT^2,\\
 B_\lambda(T)&=T(\lambda_0^Q+\beta_Q\bar\lambda_QT)e^{c_I\alpha_QT},\\
 H_J&=c_I\log c_I-c_I+1+\tfrac12c_Ib_J^2\sigma_J^2,\\
 \mathbb E^Q[Z_{T\wedge\tau_n}\log Z_{T\wedge\tau_n}]
 &\leq\gamma B_v(T)+H_JB_\lambda(T)<\infty.
\label{eq:pq-joint-entropy}
\end{align}
The jump term follows from conditional normal-mark entropy; it vanishes
when there are no jumps. Uniform stopped-density entropy implies uniform
integrability, so $\mathbb E^Q[Z_T]=1$. Nonexplosion, finite state integrals
and a strictly positive tilt imply $Z_T>0$ almost surely. Thus the two laws
are equivalent on every fixed finite horizon for this exact family.
This does not assert uniform integrability as $T\to\infty$. Stationarity is
an additional long-run restriction; a changed initial-state law,
mark-dependent excitation or state-dependent tilt needs a separate argument.
The executable entropy bound checks these coefficient conditions, not the
accuracy of a discretized likelihood or the identification of a risk premium.

GBM is the constant-variance, no-jump special case; Heston retains only the
Brownian premia; Bates adds intensity and mark premia; Hawkes also requires
the linked excitation dynamics. If a calibrated Bates model has exactly
zero jump intensity, an equivalent measure cannot create jump events on a
previously null event set. The current copper Bates fit has this boundary
case. A physical jump requirement would therefore call for a separately
versioned joint refit, not an arbitrary positive intensity inserted into the
existing fit. Very small jump marks can likewise make risk-premium estimates
weakly identified even when the algebra is admissible.

\subsection{Historical information, clocks and the commodity basis}

The additional IB API history contains 1,253 GLD regular-session daily closes
from 7 September 2021 to 2 September 2026. Its 168 overlapping observations
match the existing Paper series exactly. Full-session closes were not
interchangeable with that series and are kept separate. Daily underlying
observations serve time-series estimation; the roughly 100-contract budget
for a sparse option surface governs a different, cross-sectional sample.

Gold midpoint hourly bars and the exchange-session schedule yield 1,003
synchronous GLD--gold observations from 2 September 2022 to the cutoff,
including nine early-closing sessions. Hourly timestamps denote bar starts;
the selected bar ends at the reported equity-session close. The 250 GLD
dates without a matched gold observation are excluded, not filled. The
comparison remains one between a trade close and a midpoint from a different
market, rather than a claim of identical executions.

Let $B_t=\log(G_t/X_t)$, where $G_t$ is GLD in USD/share and $X_t$ is gold in
USD/oz. A descriptive linear fit to $B_t$ has annual slope $-0.003958$ and
residual standard deviation $0.000258$. These are empirical basis statistics,
not an identification of management fees or a permanent structural law.
An expanding-origin basis-only check gives relative conversion RMSE of
$0.0276\%$, versus $0.0358\%$ for the preceding observed ratio, over 751
targets. That check conditions on the target GLD price: it validates a
conversion diagnostic, not an ex ante gold-price forecast. A stochastic basis
model and its dependence on GLD would be needed to propagate basis risk fully.

Copper histories were obtained for seven deliveries, September 2025 through
December 2026. The requested December 2024 and March, May and July 2025
deliveries were unavailable in the attempted retrievals. The available
universe is therefore not a reconstruction of the historical front month.
Neither thin far-delivery observations nor an undocumented continuous-contract
adjustment are silently treated as spot history.

The explicit return diagnostic selects the nearest available delivery with
at least 100 contracts of volume and at least 30 calendar days to expiry,
using origin information. It then computes the next observation's return
from two prices of that same delivery. Missing or illiquid targets are
excluded without selecting a different contract retrospectively. There are
408 accepted returns from 14 November 2024 to 2 September 2026, 65 origin
failures and 14 target-liquidity exclusions. These rotating-delivery returns
avoid artificial rollover jumps but do not identify spot drift, convenience
yield, collateral returns or a five-year realised selling-price curve.

\subsection{Estimation and numerical evidence}

For GBM, actual calendar gaps enter the likelihood. If $x_j$ is a log return
over $h_j$ years, the constant log drift estimate is
$\widehat a=\sum_jx_j/\sum_jh_j$, and the arithmetic drift is
$\widehat\mu_P=\widehat a+\widehat\sigma^2/2$. On the five-year GLD sample,
the historical diffusion estimate is $20.87\%$ and arithmetic drift
$19.75\%$ annually. Holding the current option-implied diffusion at
$25.09\%$ gives drift $20.72\%$, with a model-conditional approximate
95\% interval from $-1.30\%$ to $42.74\%$. This wide interval matters for
multi-year revenues. The fixed-current-$Q$ estimate is in-sample and is not
used to score historical forecasts.

The richer preliminary fits hold the current audited Paper pricing
coefficients fixed and estimate the restricted physical premia against
historical GLD returns with a particle likelihood. They are exploratory
conditional fits, not jointly estimated historical $P/Q$ systems. In
particular, imposing today's pricing coefficients throughout a past sample
cannot support an origin-only OOS claim. The filtered physical terminal
variance and option-implied current variance must also be reconciled as
information about the same latent state.

The numerical implementation integrates latent variance and jump intensity
over substeps and weights the aggregate observed daily return. It does not
invent intraday equity observations. Independent price Brownian increments
are integrated out conditionally; state particles are resampled after each
observed return. Gaussian-limit, jump-support and coefficient-map checks
pass, but those identities are distinct from convergence of the estimated
parameters.

Refinement at 4,096 particles and three seeds changes likelihood levels:
mean negative log likelihood moves approximately from $-3874.56$ to
$-3876.82$ for Heston, from $-3875.82$ to $-3877.13$ for Bates, and from
$-3878.72$ to $-3881.11$ for Hawkes when moving from one to eight substeps
per observation. Parameters are held fixed during this check. Remaining
simulation variation and discretisation sensitivity prevent these likelihoods
from establishing a reliable model ranking. Parameter refitting at a
converged resolution, multiple starts, identification checks and uncertainty
assessment remain necessary; optimizer success alone does not satisfy them.
Further checks at 16 and 32 substeps leave seed variation: at 32 substeps
the three-seed means are approximately $-3877.65$, $-3877.41$ and
$-3880.66$, respectively. Heston and Bates even reverse their mean ordering
relative to the one-step diagnostic. This refinement does not establish
convergence of a refitted estimator or an OOS advantage for Hawkes.

\subsection{Three validation targets and the first physical baseline}

Option repricing, physical commodity prediction and operating prediction
answer different questions. The existing option test evaluates pricing
functions under $\mathbb Q$, conditional on target forwards and rates;
the observed evolution of the surface itself occurs under $\mathbb P$.
The new commodity test forecasts future prices using only origin information.
Barrick's operating test compares production, sales, costs and cash flows
with later reports. A comparison with Barrick's stock quote is not an
observed intrinsic-value target.

The first gold physical baseline estimates GBM from historical prices alone.
At each origin it separately fits the GLD--gold log-ratio trend and forecasts
the target gold price without using target GLD. It is a physical benchmark
within the GBM family, not a joint calibration to historical options. Direct
gold GBM and a no-change conditional mean, with historically estimated gold
volatility, provide comparators. Each fit uses only available origin history.
The basis requires at least 253 training prices; the GLD volatility history
is longer, so the variants differ in both input history and basis treatment.

\begin{table}[H]
\centering\small
\caption{Training-only physical gold baselines, cutoff 2 September 2026.
RMSE is in USD/oz; coverage is realised coverage of nominal 95\% plug-in
intervals. Quarterly forecasts overlap.}
\label{tab:physical-gold-baseline}
\begin{tabularx}{\textwidth}{@{}Xrrrr@{}}
\toprule
 & \multicolumn{2}{c}{Next session: 750 targets}
 & \multicolumn{2}{c}{90 days: 689 targets}\\
Model & RMSE & Coverage & RMSE & Coverage\\
\midrule
GLD GBM plus basis trend & 53.82 & 89.3\% & 417.19 & 89.6\%\\
Direct gold GBM & 53.77 & 90.0\% & 421.17 & 92.6\%\\
Gold no-change mean & 53.86 & 89.7\% & 428.52 & 82.4\%\\
\bottomrule
\end{tabularx}
\end{table}

Quarterly targets are the first observed session at or after 90 calendar
days. They are not 689 independent experiments. The GLD/basis variant has
slightly lower quarterly RMSE but worse absolute error and average log score
than direct gold GBM; there is no unanimous ranking. All reported 95\%
coverages fall below their nominal level. Parameter uncertainty and stochastic
basis uncertainty are omitted from these plug-in intervals. The results
therefore establish an executable physical evaluation, not adequate long-run
distribution calibration or superiority of the option-calibrated engines.

For the richer models, every historical origin must use the pricing
calibration and underlying history available at that origin, with latent-state
filtering and physical premium estimation completed before the target is
observed. Daily distribution scores, interval coverage and quarterly price
errors must remain separate from option IV errors. Reusing a September
calibration to forecast earlier dates would violate this design.

Thirty historical surfaces have now been recovered from the observed fields
of the fixed rolling panel, each with 64 contracts across at least three
expiries. The recovery excludes prior model predictions and support-selection
flags; all dated spot levels agree with independently acquired GLD regular-session
closes. These are observed legacy option-price fields, not asserted transaction
prices. Recovering the observations does not recover historical model parameters:
each origin still requires its own calibration. The window supports short-horizon
experiments, but no 90-day target after the first July origin lies within the
2 September market cutoff.

Joint maximum likelihood is not a mathematical requirement of changing measure.
A conditional two-stage estimator can calibrate structural pricing coefficients
on origin-date options and then estimate physical premia from returns available
by that origin, holding those coefficients fixed. This imposes their constancy
over the estimation history and must not be described as joint maximum likelihood.
The forward simulation conditions on the option-estimated current variance,
shared under both measures; the historical return filter starts from its stated
physical long-run mean. Its terminal estimate provides a separate diagnostic,
not an automatic replacement for the option state. Parameter and state-estimation
uncertainty from both stages must be assessed before interpreting plug-in
intervals as unconditional predictive coverage.

The completed conditional GBM experiment has 29 next-session targets, from
10 July--1 September origins. Each origin has a new option-volatility fit;
the physical log drift is estimated from returns available by that origin.
Exact lognormal prediction gives RMSE of USD~6.119 per GLD share against
USD~6.117 for the no-change mean. In the matched four-model panel below,
\CovBSNinetyFive{} nominal 95\% gold intervals cover the observed price. This short sample supplies no RMSE
advantage, and it does not establish long-horizon calibration. These GLD
errors are distinct from the USD/oz errors of the longer historical-only
gold benchmark above. The richer-model origin experiment must additionally
resolve particle and time-discretisation sensitivity.

The Heston panel has subsequently completed the same 29 forecast origins.
Its GLD RMSE is USD~6.157 per share, against 6.119 for conditional GBM and
6.117 for no-change. Two physical optimisations exhaust their evaluation
budget; their forecasts remain in the comparison rather than being removed
after observing the errors. Applying the origin-fitted deterministic gold
basis to both models gives RMSE of USD~66.77/oz for Heston and 66.35/oz for
GBM, against 66.09/oz for gold no-change. Both nominal 95\% intervals cover
\CovHNinetyFive{} gold targets. These matched-sample results establish no RMSE
advantage, and do not certify parameter identification or long-horizon
distributional accuracy.

The Bates and Bates--Hawkes panels subsequently completed all 29 matched
origins, including the separately saved 10 July pilot. The common targets
and the origin-only gold conversion exclude no observations from any model.
Table~\ref{tab:physical-four-family-complete} supersedes the single-origin
comparison as the descriptive panel result. Bates--Hawkes has a slightly
lower RMSE than no-change on both scales; this short sample does not
establish statistical superiority or long-horizon accuracy. Its gold MAE
is USD~51.98/oz, whereas Bates has USD~51.46/oz, so the apparent ranking
also depends on the loss function.

\begin{table}[H]
\centering\small
\caption{Completed conditional physical experiment: 29 common next-session
targets, 13 July--2 September 2026. Gold conversion uses an origin-estimated
deterministic basis; parameter and basis uncertainty are omitted.}
\label{tab:physical-four-family-complete}
\begin{tabular}{lrrcc}
\toprule
Model & GLD RMSE & Gold RMSE & Gold 95\% & Gold 80\%\\
 & USD/share & USD/oz & coverage & coverage\\
\midrule
GBM & 6.119 & 66.35 & \CovBSNinetyFive & \CovBSEighty\\
Heston & 6.157 & 66.77 & \CovHNinetyFive & \CovHEighty\\
Bates & 6.162 & 66.84 & \CovBNinetyFive & \CovBEighty\\
Bates--Hawkes & 6.072 & 65.81 & \CovWNinetyFive & \CovWEighty\\
No-change mean & 6.117 & 66.09 & --- & ---\\
\bottomrule
\end{tabular}
\end{table}

Coverage is tested against the binomial law of the nominal level. At 95\%
the jump models cover \CovBNinetyFive{} targets (binomial $p=\CovBNinetyFiveP$),
against \CovBSNinetyFive{} for GBM and Heston ($p=\CovBSNinetyFiveP$). At 80\%
the jump models under-cover significantly: Bates covers \CovBEighty{}
($p=\CovBEightyP$) and Bates--Hawkes \CovWEighty{} ($p=\CovWEightyP$), against
\CovBSEighty{} for GBM ($p=\CovBSEightyP$) and \CovHEighty{} for Heston
($p=\CovHEightyP$). The central intervals of the jump models are too narrow,
and parameter and basis uncertainty are omitted. These origin fits predate the
24 October carry correction: their pricing stage uses $q=0$, and they were not
recalibrated with $q=0.40\%$.

\paragraph{Numerical reliability and chronological initialization.}
Two Heston physical optimizations (24 July and 20 August) and one
Bates--Hawkes optimization (19 August) report failure; their forecasts
remain included. Across three independent seeds at each of two time grids,
the fitted physical premia improve on zero premia in every check at only
10/29 Heston, 6/29 Bates and 9/29 Bates--Hawkes origins. The maximum
within-grid negative-log-likelihood spread across seeds is approximately
6.08, 455.38 and 402.10 respectively. These are numerical diagnostics,
not confidence intervals: optimizer termination does not certify a stable
likelihood estimate or identified premia.

The resumed computation initializes from the latest successful strictly
earlier origin, then refits on the current origin's information. Heston
and Bates first try constrained local refinement with global fallback;
Bates--Hawkes also retains a current-origin nested Bates candidate. The
physical stage compares earlier premia with zero premia on the current
training likelihood before choosing its start. Saved artifacts record
the source date and hash. Existing cold-start results are preserved, so
this completed panel has a documented mixture of initialization policies;
it is not a uniform warm-start experiment. A same-surface numerical check
found different local minima for Heston and Bates, which prevents claiming
that the faster initialization reproduces the original global optimum.
The artifact audit verifies 58 pricing-to-physical hash links and 49
strictly backward initialization links. Those checks establish traceability,
not numerical convergence. The aggregate output versions are
\texttt{20261011-four-model-complete} and
\texttt{20261011-four-model-gold-complete}.

The original 10 July pilot, trained on 1,214 GLD returns and forecasting
13 July, is included in the panel. Its gold absolute errors were
USD~114.04, 119.31, 131.91 and 105.07 per ounce for GBM, Heston, Bates
and Bates--Hawkes respectively,
against 106.45 for no-change. Conversion fits the log GLD/gold ratio only
through each origin and multiplies the forecast mean and quantiles by the
predicted factor, without using the target GLD price. The gold no-change
benchmark uses the actual origin gold price. Gold midpoints remain distinct
from Barrick's realised selling prices.

\paragraph{Analytical integration of normal marks.}
An optional filter variant integrates jump amplitudes analytically. Given
the simulated variance path and event counts, a total of $N$ independent
normal marks contributes mean $N\mu_{J,P}$ and variance $N\sigma_J^2$ to
the conditional daily-return density. This is valid for the implemented
nonmark-dependent intensity and variance dynamics. Event counts and
variance paths remain sampled; the fixed panel above retains its original
filter. The new variant matches the one-step compound-Poisson normal-mixture
reference and leaves the zero-jump filter unchanged.

A targeted check uses 24 July, selected separately for each jump family
because it has the largest saved within-grid seed spread. Holding premia
fixed, analytical marks reduce the maximum spread across the two tested
grids from 455.38 to 152.97 for Bates and from 402.10 to 139.39 for Hawkes.
Substantial sensitivity remains, including a larger Hawkes spread at the
finer grid. This is a numerical stress check, not a refit, new OOS evidence
or a convergence certificate; it identifies further particle/count
integration work before relying on long-horizon physical distributions.

\subsection{Consequences for the operating and valuation layers}

The required revenue input is $\mathbb E^P[P_t^{\rm realised}Q_t^{\rm sold}]$
on the accepted ownership basis. It reduces to a deterministic sales volume
times the expected realised price only under that operating simplification;
with stochastic sales, dependence matters. Royalties, price-linked taxes,
inventory timing and realised-price basis also belong in the operating
transformation. Depreciation inside Cost of Sales must be reconciled with
cash capital expenditure and the EBIT-to-FCFF bridge.

Instrument identity and the valuation clock impose further restrictions. A
fixed-delivery copper future cannot be continued beyond expiry as a spot
process to 2030: rotating-delivery returns alone do not identify that extension.
A copper term-structure and realised-price basis are required. Likewise, a
2 September valuation excludes already elapsed first- and second-quarter cash
flows from future enterprise value. The remaining third-quarter stub, later
quarters and balance-sheet claims must use a consistent valuation date.

The fixed option-shape scenario law remains a separate experiment. None of
its valuation quantiles is retroactively relabelled physical because a
coefficient map or a preliminary estimate now exists. A replacement requires
admissible and identified physical models, origin-only validation, uncertainty
propagation, compatible discounting and reconciled corporate inputs. The
following chapters supply the operating quantities and expose the additional
accounting conditions that must hold before that replacement can be made.
